\documentclass[conference]{IEEEtran}
\IEEEoverridecommandlockouts

% Packages
\usepackage{cite}
\usepackage{amsmath,amssymb,amsfonts}
\usepackage{algorithmic}
\usepackage{graphicx}
\usepackage{textcomp}
\usepackage{xcolor}
\usepackage{booktabs}
\usepackage{tabularx}
\usepackage{multirow}
\usepackage{hyperref}
\usepackage{listings}
\usepackage{float}

% Code snippet formatting
\lstset{
    backgroundcolor=\color{gray!10},
    basicstyle=\ttfamily\footnotesize,
    breaklines=true,
    captionpos=b,
    commentstyle=\color{green!60!black},
    keywordstyle=\color{blue},
    stringstyle=\color{red},
    frame=single
}

\begin{document}

\title{LUMEN: Security and Quality Evaluation of LLM-Generated Code}

\author{
\IEEEauthorblockN{
    Hafiz Muhammad Hassam Ata, 
    Muhammad Ayan Tahir, 
    Muhammad Anas Khan, 
    Abdul Ahad Khan
}
\vspace{0.15cm}
\IEEEauthorblockA{
    \textit{Department of Software Engineering} \\
    \textit{NED University of Engineering and Technology}\\
    Karachi, Pakistan \\
    \{ata4501984, tahir4506733, khan4501430, khan4530246\}@cloud.neduet.edu.pk
}
}

\maketitle

\begin{abstract}
% TODO: Revise abstract to reflect your current research findings if needed.
Although the large language models (LLMs) used by many developers are gaining popularity, and software development productivity can be enhanced by them, it is worth investigating their security and quality impact on the code they are generating. This helps avoid the introduction of high vulnerabilities and maintenance problems to software systems over time. Nevertheless, the latest studies on automated code generation are narrow and lack in giving an exhaustive review. Systems with a focus on both security and functionality are rare; most common benchmarks mainly test functional correctness, and ignore hidden quality concerns and vulnerabilities. The primary aim of this paper is to test the capacity of different state-of-the-art LLMs---to be precise ChatGPT 5.3 mini, Gemini 3 Flash, and Grok 4 fast---to create safe, high-quality and testable code only in the JavaScript language, benchmarked in direct opposition to man-made baselines. The collected data were through manually generated datasets of zero-shot, functionally-driven prompts in a variety of five risk zones, and were simulated with Lumen, a hybrid analysis platform which combines deterministic static analysis through SonarQube with a smart semantic assessment system based on a purpose-trained Qwen 3 4B model (Q4 GGUF), fine-tuned on the CrossVul vulnerability dataset \cite{b25} to output structured SARIF 2.1.0 findings. The fine-tuned judge achieves 94\% accuracy with perfect recall on a held-out balanced evaluation set, demonstrating the viability of a lightweight, locally-deployable SAST model as an effective triage engine. The outcome-based results show that insecure code of high-severity problems is commonly generated by LLMs, such as arbitrary code execution by unsafe deserialization, SSRF, and hard coded cryptographic salts. In addition, the notion of functional correctness is not a reliable indicator of overall code security, which is successful code that often passes unit tests often contains inherent weaknesses due to uncritical action on blind obedience. Furthermore, baselines written by humans always scored higher than AI models through the use of defensive program constructs, dramatic input sanitization, and environment isolation. These results imply that developers need to use stringent validation practices, and thoroughly test and investigate any generated code with respect to security before deployment. Human oversight and hybrid assessment tools such as Lumen should be applied in organizations as vital elements of safe software development and protection of the supply chain.
\end{abstract}

\begin{IEEEkeywords}
Large Language Models, Code Security, Automated Code Generation, Zero-Shot Prompting, Vulnerability Detection, SonarQube, Qwen 3, Fine-tuning, CrossVul, SARIF, Machine Learning, JavaScript
\end{IEEEkeywords}

\section{Introduction}
% TODO: Write your introduction context here. Introduce the problem domain, motivation, and objectives.

\subsection{Background and Motivation}
% TODO: Add background details and problem statement.

\subsection{Proposed Approach (Lumen Framework)}
% TODO: Explain the hybrid platform architecture (SonarQube + Qwen 3).

% TODO: PICTURE PLACEHOLDER - Add an overview architecture diagram of Lumen here.
% \begin{figure}[htbp]
%     \centering
%     \includegraphics[width=\linewidth]{lumen_architecture.png}
%     \caption{Overview of the Lumen hybrid evaluation architecture.}
%     \label{fig:lumen_architecture}
% \end{figure}

\subsection{Research Questions}
To guide this study, we formulated the following Research Questions (RQs):
\begin{itemize}
    \item \textbf{RQ1:} How secure is the code generated by LLMs compared to human-written code?
    \item \textbf{RQ2:} What are the most common vulnerabilities found in LLM-generated code?
    \item \textbf{RQ3:} How effective is the hybrid Lumen platform in evaluating code security?
\end{itemize}

\section{Literature Review}

\subsection{Introduction}
% TODO: Add introductory literature context.

\subsection{Security Risks in LLM-Generated Code}
% TODO: Write literature discussion on LLM security risks and vulnerability taxonomies.

% TODO: PICTURE PLACEHOLDER - Vulnerability distribution chart (Figure 1).
\begin{figure*}[tbp]
    \centering
    % \includegraphics[width=0.65\linewidth]{figure1.png}
    \caption{Distribution of security vulnerability categories identified in LLM-generated code.}
    \label{fig:vulnerability_distribution}
\end{figure*}

\subsection{Code Evaluation Static Analysis Tools}
% TODO: Discuss SAST tools, limitations, and tool agreement metrics.

% TODO: PICTURE PLACEHOLDER - SAST Venn Diagram.
\begin{figure*}[tbp]
\centering
% \includegraphics[width=0.6\linewidth]{figure_sast_venn.png}
\caption{Venn diagram displaying the low overlap of detected Vulnerability-Contributing Commits among different static analysis tools.}
\label{fig:venn_diagram}
\end{figure*}

% TODO: PICTURE PLACEHOLDER - Root causes distribution chart.
\begin{figure*}[tbp]
\centering
% \includegraphics[width=1\linewidth]{figure_root_causes.png}
\caption{Distribution of input characteristics leading to False Negatives and False Positives in static code analyzers.}
\label{fig:root_causes}
\end{figure*}

\subsection{AI Based Code Quality and Vulnerability Detection}
% TODO: Discuss machine learning based classification and LLM-as-a-judge approaches.

% TODO: PICTURE PLACEHOLDER - SAST + AI Genius Pipeline.
\begin{figure*}[tbp]
    \centering
    % \includegraphics[width=0.8\linewidth]{figure_sast_genius.png} 
    \caption{A hybrid vulnerability detection framework where deterministic static analysis findings are evaluated and filtered by an intelligent AI triage engine.}
    \label{fig:sast_genius}
\end{figure*}

% TODO: PICTURE PLACEHOLDER - Iris pipeline diagram.
\begin{figure*}[tbp]
    \centering
    % \includegraphics[width=1\linewidth]{figure_iris_pipeline.png} 
    \caption{A pipeline demonstrating how language models can infer rules for unknown libraries and feed them into a static analysis engine.}
    \label{fig:iris_pipeline}
\end{figure*}

\subsection{Comparison Between Human and AI-Written Code}
% TODO: Detail functional correctness vs. defensive programming between human and AI implementations.

\section{Methodology}
% TODO: Detail your experimental setup, prompt datasets, risk zones, evaluation metrics, and configuration of models (ChatGPT 5.3 mini, Gemini 3 Flash, Grok 4 fast).

% TODO: PICTURE PLACEHOLDER - Add experimental pipeline or dataset generation flow chart here.

\section{Results and Evaluation}
% TODO: Present findings answering RQ1, RQ2, and RQ3 with tables and charts.

% TODO: TABLE PLACEHOLDER - Add comparison table of model performance vs human baselines.

\section{Discussion}
% TODO: Discuss implications of the findings, threats to validity, and limitations.

\section{Conclusion and Future Work}
% TODO: Summarize key contributions and outline future research directions.

\begin{thebibliography}{00}
% TODO: Add your references here (e.g., b1, b8, b25).
\bibitem{b1} Author, ``Title'', Journal, Year.
\bibitem{b8} Author, ``Title'', Conference, Year.
\bibitem{b25} Author, ``CrossVul Dataset'', Year.
\end{thebibliography}

\end{document}